\begin{abstract}
  The rise of machine learning has shifted targeted resource allocation in social and humanitarian settings toward algorithmic targeting based on predicted risk scores. This approach is typically cheaper and faster than traditional screening procedures that directly observe the latent vulnerability status through household surveys or technical assessments. Yet, even access to the true conditional vulnerability probability cannot eliminate misallocation: aleatoric uncertainty over individual vulnerability status persists irreducibly, and probabilistic targeting inevitably misallocates some resources. In this work we study how screening and algorithmic targeting should be optimally combined in a two-stage allocation framework where a screening stage observes true outcomes for a subset of units before a final allocation stage assigns the resource under a fixed coverage budget. We show that the optimal screening set consists of units "in the middle" of the risk distribution, where uncertainty most affects allocation decisions, while directly targeting the highest-risk units. We derive the optimal screening set in closed form for uniform risk and provide a fixed-point algorithm that recovers it for arbitrary risk distributions to operationalize our results. Furthermore, we empirically characterize when screening and algorithmic targeting act as complements or substitutes: efficiency gains from screening grow as the aleatoric uncertainty in the population increases. We illustrate our framework in humanitarian demining and poverty alleviation programs, where the tension between screening costs and allocation efficiency is operationally consequential.
\end{abstract}